\documentclass[10pt,a4paper]{amsart}
\usepackage{amsmath,amssymb,amsthm,graphicx}
% For dealing with references we use the comment environment
\usepackage{verbatim}
\newenvironment{reference}{\comment}{\endcomment}
%\newenvironment{reference}{}{}
\newenvironment{slogan}{\comment}{\endcomment}
\newenvironment{history}{\comment}{\endcomment}

% For commutative diagrams we use Xy-pic
\usepackage[all]{xy}
% Theorem environments.
%
\theoremstyle{plain}
\newtheorem{theorem}[subsection]{Theorem}
\newtheorem{proposition}[subsection]{Proposition}
\newtheorem{lemma}[subsection]{Lemma}

\theoremstyle{definition}
\newtheorem{definition}[subsection]{Definition}
\newtheorem{example}[subsection]{Example}
\newtheorem{exercise}[subsection]{Exercise}
\newtheorem{situation}[subsection]{Situation}

\theoremstyle{remark}
\newtheorem{remark}[subsection]{Remark}
\newtheorem{remarks}[subsection]{Remarks}

\numberwithin{equation}{subsection}

% Macros
%
\def\lim{\mathop{\mathrm{lim}}\nolimits}
\def\colim{\mathop{\mathrm{colim}}\nolimits}
\def\Spec{\mathop{\mathrm{Spec}}}
\def\Hom{\mathop{\mathrm{Hom}}\nolimits}
\def\Ext{\mathop{\mathrm{Ext}}\nolimits}
\def\SheafHom{\mathop{\mathcal{H}\!\mathit{om}}\nolimits}
\def\SheafExt{\mathop{\mathcal{E}\!\mathit{xt}}\nolimits}
\def\Sch{\mathit{Sch}}
\def\Mor{\mathop{\mathrm{Mor}}\nolimits}
\def\Ob{\mathop{\mathrm{Ob}}\nolimits}
\def\Sh{\mathop{\mathit{Sh}}\nolimits}
\def\NL{\mathop{N\!L}\nolimits}
\def\CH{\mathop{\mathrm{CH}}\nolimits}
\def\proetale{{pro\text{-}\acute{e}tale}}
\def\etale{{\acute{e}tale}}
\def\QCoh{\mathit{QCoh}}
\def\Ker{\mathop{\mathrm{Ker}}}
\def\Im{\mathop{\mathrm{Im}}}
\def\Coker{\mathop{\mathrm{Coker}}}
\def\Coim{\mathop{\mathrm{Coim}}}

% Boxtimes
%
\DeclareMathSymbol{\boxtimes}{\mathbin}{AMSa}{"02}

%
% Macros for moduli stacks/spaces
%
\def\QCohstack{\mathcal{QC}\!\mathit{oh}}
\def\Cohstack{\mathcal{C}\!\mathit{oh}}
\usepackage{cleveref}
\crefname{originalproposition}{original proposition}{original propositions}
\Crefname{originalproposition}{Original proposition}{Original propositions}
\makeatletter
\newlabel{proposition-homotopy-category-triangulated}{{014S}{source}{}{sourcepointer-1}{}}
\newlabel{proposition-homotopy-category-triangulated@cref}{{[originalproposition][0][]014S}{[1][0][]source}{}{}{}}
\makeatother
\crefname{originallemma}{original lemma}{original lemmas}
\Crefname{originallemma}{Original lemma}{Original lemmas}
\makeatletter
\newlabel{lemma-subcategory-left-resolution}{{05T7}{source}{}{sourcepointer-2}{}}
\newlabel{lemma-subcategory-left-resolution@cref}{{[originallemma][0][]05T7}{[1][0][]source}{}{}{}}
\makeatother
\crefname{originallemma}{original lemma}{original lemmas}
\Crefname{originallemma}{Original lemma}{Original lemmas}
\makeatletter
\newlabel{lemma-improve-distinguished-triangle-homotopy}{{0G6C}{source}{}{sourcepointer-3}{}}
\newlabel{lemma-improve-distinguished-triangle-homotopy@cref}{{[originallemma][0][]0G6C}{[1][0][]source}{}{}{}}
\makeatother
\title{Compatible bounded-above resolutions with split transition maps for an arbitrary complex}
\author{The Stacks Project Authors}
\date{stacks-project:derived.tex:local-snapshot-20261011; GNU FDL 1.2 or later}
\begin{document}
\maketitle
\noindent\textit{Editorial note.} P is only closed under finite direct sums and epimorphism-generating; its members need not be projective. The split-injection homotopy adjustment and final contractible correction remain counted work.\par\medskip
\noindent\textit{Editorial scope.} One original target with its local construction proved in full. Classical category primitives are explicitly allowed.\par\medskip
\noindent\textit{Allowed standard prerequisites (editorial).}\begin{itemize}
\item Bounded-above resolutions by an epimorphism-generating class. A bounded-above complex has a termwise surjective quasi-isomorphism from a bounded-above complex in P if every object is a quotient of an object in P.
\item Homotopy-category triangles. The homotopy category of an abelian category is triangulated; any distinguished triangle is isomorphic to one arising from a termwise split exact sequence. Standard cone formulas and long exact cohomology are allowed.
\end{itemize}\medskip
\ifdefined\hypertarget\hypertarget{sourcepointer-1}{}\fi
\noindent\textit{Source reference.} Original proposition 014S: focused-seeds/ACL/sources/stacks-snapshot-20261011/derived.tex, source lines 3319--3326. Standard structure of the homotopy category. Stable Stacks tag displayed.\par\medskip
\ifdefined\hypertarget\hypertarget{sourcepointer-2}{}\fi
\noindent\textit{Source reference.} Original lemma 05T7: focused-seeds/ACL/sources/stacks-snapshot-20261011/derived.tex, source lines 5301--5316. Explicit permitted single-complex bounded-above resolution theorem; compatibility and split tower repair are not granted. Stable Stacks tag displayed.\par\medskip
\ifdefined\hypertarget\hypertarget{sourcepointer-3}{}\fi
\noindent\textit{Source reference.} Original lemma 0G6C: focused-seeds/ACL/sources/stacks-snapshot-20261011/derived.tex, source lines 3458--3466. Standard split-cone realization of a homotopy triangle. Stable Stacks tag displayed.\par\medskip
\begin{lemma}
\label{lemma-special-direct-system}
Let $\mathcal{A}$ be an abelian category. Let
$\mathcal{P} \subset \Ob(\mathcal{A})$ be a subset.
Assume $\mathcal{P}$ contains $0$, is closed under (finite) direct sums,
and every object of $\mathcal{A}$ is a quotient of an
element of $\mathcal{P}$. Let $K^\bullet$ be a complex.
There exists a commutative diagram
$$
\xymatrix{
P_1^\bullet \ar[d] \ar[r] & P_2^\bullet \ar[d] \ar[r] & \ldots \\
\tau_{\leq 1}K^\bullet \ar[r] & \tau_{\leq 2}K^\bullet \ar[r] & \ldots
}
$$
in the category of complexes such that
\begin{enumerate}
\item the vertical arrows are quasi-isomorphisms and termwise surjective,
\item $P_n^\bullet$ is a bounded above complex with terms in
$\mathcal{P}$,
\item the arrows $P_n^\bullet \to P_{n + 1}^\bullet$
are termwise split injections and each cokernel
$P^i_{n + 1}/P^i_n$ is an element of $\mathcal{P}$.
\end{enumerate}
\end{lemma}

\begin{proof}
We are going to use that the homotopy category $K(\mathcal{A})$ is a
triangulated category, see Proposition
\ref{proposition-homotopy-category-triangulated}.
By Lemma \ref{lemma-subcategory-left-resolution} we can find a
termwise surjective map of complexes $P_1^\bullet \to \tau_{\leq 1}K^\bullet$
which is a quasi-isomorphism
such that the terms of $P_1^\bullet$ are in $\mathcal{P}$.
By induction it suffices, given
$P_1^\bullet, \ldots, P_n^\bullet$ to construct
$P_{n + 1}^\bullet$ and the maps
$P_n^\bullet \to P_{n + 1}^\bullet$ and
$P_{n + 1}^\bullet \to \tau_{\leq n + 1}K^\bullet$.

\medskip\noindent
Choose a distinguished triangle
$P_n^\bullet \to \tau_{\leq n + 1}K^\bullet \to C^\bullet \to P_n^\bullet[1]$
in $K(\mathcal{A})$. Applying
Lemma \ref{lemma-subcategory-left-resolution} we choose a
map of complexes $Q^\bullet \to C^\bullet$
which is a quasi-isomorphism such that the terms of $Q^\bullet$
are in $\mathcal{P}$. By the axioms of triangulated categories
we may fit the composition $Q^\bullet \to C^\bullet \to P_n^\bullet[1]$ into
a distinguished triangle
$P_n^\bullet \to P_{n + 1}^\bullet \to Q^\bullet \to P_n^\bullet[1]$
in $K(\mathcal{A})$.
By Lemma \ref{lemma-improve-distinguished-triangle-homotopy}
we may and do assume
$0 \to P_n^\bullet \to P_{n + 1}^\bullet \to Q^\bullet \to 0$
is a termwise split short exact sequence. This implies that
the terms of $P_{n + 1}^\bullet$ are in $\mathcal{P}$ and that
$P_n^\bullet \to P_{n + 1}^\bullet$ is a termwise split injection
whose cokernels are in $\mathcal{P}$.
By the axioms of triangulated categories we obtain a map
of distinguished triangles
$$
\xymatrix{
P_n^\bullet \ar[r] \ar[d] &
P_{n + 1}^\bullet \ar[r] \ar[d] &
Q^\bullet \ar[r] \ar[d] &
P_n^\bullet[1] \ar[d] \\
P_n^\bullet \ar[r] &
\tau_{\leq n + 1}K^\bullet \ar[r] &
C^\bullet \ar[r] &
P_n^\bullet[1]
}
$$
in the triangulated category $K(\mathcal{A})$. Choose an actual morphism of
complexes $f : P_{n + 1}^\bullet \to \tau_{\leq n + 1}K^\bullet$.
The left square of the diagram above commutes up to homotopy, but as
$P_n^\bullet \to P_{n + 1}^\bullet$ is a termwise split injection
we can lift the homotopy and modify our choice of $f$ to make it commute.
Finally, $f$ is a quasi-isomorphism, because both $P_n^\bullet \to P_n^\bullet$
and $Q^\bullet \to C^\bullet$ are.

\medskip\noindent
At this point we have all the properties we want, except we don't know
that the map $f : P_{n + 1}^\bullet \to \tau_{\leq n + 1}K^\bullet$
is termwise surjective. Since we have the commutative diagram
$$
\xymatrix{
P_n^\bullet \ar[d] \ar[r] & P_{n + 1}^\bullet \ar[d] \\
\tau_{\leq n}K^\bullet \ar[r] & \tau_{\leq n + 1}K^\bullet
}
$$
of complexes, by induction hypothesis we see that $f$ is surjective
on terms in all degrees except possibly $n$ and $n + 1$. Choose
an object $P \in \mathcal{P}$ and a surjection $q : P \to K^n$.
Consider the map
$$
g :
P^\bullet = (\ldots \to 0 \to P \xrightarrow{1} P \to 0 \to \ldots)
\longrightarrow
\tau_{\leq n + 1}K^\bullet
$$
with first copy of $P$ in degree $n$ and maps given by
$q$ in degree $n$ and $d_K \circ q$ in degree $n + 1$.
This is a surjection in degree $n$ and the cokernel in
degree $n + 1$ is $H^{n + 1}(\tau_{\leq n + 1}K^\bullet)$;
to see this recall that $\tau_{\leq n + 1}K^\bullet$ has
$\Ker(d_K^{n + 1})$ in degree $n + 1$.
However, since $f$ is a quasi-isomorphism we know that
$H^{n + 1}(f)$ is surjective. Hence after replacing 
$f : P_{n + 1}^\bullet \to \tau_{\leq n + 1}K^\bullet$
by
$f \oplus g : P_{n + 1}^\bullet \oplus P^\bullet \to \tau_{\leq n + 1}K^\bullet$
we win.
\end{proof}

\end{document}
